% GENERATED FILE. Edit canonical lessons, then run npm run generate:books.
% canonical-source-hash: fnv1a64:93b646314009c44d
% canonical-lessons: BN-C377-unnoti, BN-C377-ogrogoti, BN-C377-goti, BN-C377-otikrom-kora, BN-C377-sadhon-kora

\chapter{Getting There}
\label{ch:bn-getting-there}
\hlchaptermodality{377}

\begin{chapteropening}
\textbf{By the end of this chapter:} \emph{I can say how much progress has been made and what is still ahead.}

Complete the last lesson of chapter 377: I can say how much progress has been made and what is still ahead.
\end{chapteropening}

\section[\texorpdfstring{\bn{উন্নতি} (unnôti)}{উন্নতি (unnôti)}]{\bn{উন্নতি} (unnôti) — improvement; progress}

\label{lesson:BN-C377-unnoti}



\begin{quote}
Before the new one: say the Bengali for successful, then the Bengali for to achieve.
\end{quote}

\subsection*{You'll want to know: \bn{উন্নতি}}
\textbf{\bn{উন্নতি}} — \emph{unnôti} — \textquotedblleft{}improvement; progress\textquotedblright{}.

\begin{quote}
\bn{রোজ} \bn{পড়লে} \bn{তোমার} \bn{বাংলার} \bn{উন্নতি} \bn{হবে।}
\end{quote}

\begin{quote}
\emph{roj pôṛle tomār bānglār unnôti hôbe.}
\end{quote}

\begin{quote}
If you read every day, your Bengali will improve.
\end{quote}

\begin{cousinweb}
\textbf{\bn{উন্নতি}} (\emph{unnôti}) is Sanskrit \emph{unnati}, a rising up (Monier-Williams). Bengali uses it for any getting better: a skill, a town, a patient.
\end{cousinweb}

\begin{grammarlens}[title={What you've built}]
One more word for saying what happens if.
\end{grammarlens}

\subsection*{Your turn}
\begin{itemize}
\raggedright
  \item \emph{Say it:} \emph{unnôti}
  \item \emph{Say it:} \emph{unnôti}, once more
  \item \emph{Say it:} \emph{roj pôṛle tomār bānglār unnôti hôbe.} — If you read every day, your Bengali will improve.
  \item \emph{Recall:} say the Bengali for successful, then the Bengali for to achieve, then say \emph{unnôti} again
\end{itemize}

\begin{tcolorbox}[breakable,colback=teal!4,colframe=teal!35!black,title={Before you move on}]
What does \bn{উন্নতি} mean? (\textquotedblleft{}Improvement; progress\textquotedblright{}.) Say it once more.
\end{tcolorbox}
\section[\texorpdfstring{\bn{অগ্রগতি} (ôgrôgôti)}{অগ্রগতি (ôgrôgôti)}]{\bn{অগ্রগতি} (ôgrôgôti) — progress (towards a goal)}

\label{lesson:BN-C377-ogrogoti}



\begin{quote}
Before the new one: say the Bengali for to achieve, then the Bengali for improvement.
\end{quote}

\subsection*{You'll want to know: \bn{অগ্রগতি}}
\textbf{\bn{অগ্রগতি}} — \emph{ôgrôgôti} — \textquotedblleft{}progress (towards a goal)\textquotedblright{}.

\begin{quote}
\bn{কাজের} \bn{অগ্রগতি} \bn{ভালো} \bn{হলে} \bn{আমরা} \bn{আগেই} \bn{শেষ} \bn{করব।}
\end{quote}

\begin{quote}
\emph{kājer ôgrôgôti bhālo hôle āmrā āgei shesh kôrbô.}
\end{quote}

\begin{quote}
If the work makes good progress, we'll finish early.
\end{quote}

\begin{cousinweb}
\textbf{\bn{অগ্রগতি}} (\emph{ôgrôgôti}) is Sanskrit \emph{agra}, the front, with \emph{gati}, going (Monier-Williams): going forward. It measures how far a job has come.
\end{cousinweb}

\begin{grammarlens}[title={What you've built}]
One more word for saying what happens if.
\end{grammarlens}

\subsection*{Your turn}
\begin{itemize}
\raggedright
  \item \emph{Say it:} \emph{ôgrôgôti}
  \item \emph{Say it:} \emph{ôgrôgôti}, once more
  \item \emph{Say it:} \emph{kājer ôgrôgôti bhālo hôle āmrā āgei shesh kôrbô.} — If the work makes good progress, we'll finish early.
  \item \emph{Recall:} say the Bengali for to achieve, then the Bengali for improvement, then say \emph{ôgrôgôti} again
\end{itemize}

\begin{tcolorbox}[breakable,colback=teal!4,colframe=teal!35!black,title={Before you move on}]
What does \bn{অগ্রগতি} mean? (\textquotedblleft{}Progress (towards a goal)\textquotedblright{}.) Say it once more.
\end{tcolorbox}
\section[\texorpdfstring{\bn{গতি} (gôti)}{গতি (gôti)}]{\bn{গতি} (gôti) — speed; pace}

\label{lesson:BN-C377-goti}



\begin{quote}
Before the new one: say the Bengali for improvement, then the Bengali for progress.
\end{quote}

\subsection*{You'll want to know: \bn{গতি}}
\textbf{\bn{গতি}} — \emph{gôti} — \textquotedblleft{}speed; pace\textquotedblright{}.

\begin{quote}
\bn{গতি} \bn{না} \bn{কমালে} \bn{দুর্ঘটনা} \bn{হতে} \bn{পারে।}
\end{quote}

\begin{quote}
\emph{gôti nā kômāle durghôṭônā hôte pāre.}
\end{quote}

\begin{quote}
If you don't slow down, there could be an accident.
\end{quote}

\begin{cousinweb}
\textbf{\bn{গতি}} (\emph{gôti}) is Sanskrit \emph{gati}, going or motion (Monier-Williams): how fast something goes.
\end{cousinweb}

\begin{grammarlens}[title={What you've built}]
One more word for saying what happens if.
\end{grammarlens}

\subsection*{Your turn}
\begin{itemize}
\raggedright
  \item \emph{Say it:} \emph{gôti}
  \item \emph{Say it:} \emph{gôti}, once more
  \item \emph{Say it:} \emph{gôti nā kômāle durghôṭônā hôte pāre.} — If you don't slow down, there could be an accident.
  \item \emph{Recall:} say the Bengali for improvement, then the Bengali for progress, then say \emph{gôti} again
\end{itemize}

\begin{tcolorbox}[breakable,colback=teal!4,colframe=teal!35!black,title={Before you move on}]
What does \bn{গতি} mean? (\textquotedblleft{}Speed; pace\textquotedblright{}.) Say it once more.
\end{tcolorbox}
\section[\texorpdfstring{\bn{অতিক্রম} \bn{করা} (ôtikrôm kôrā)}{অতিক্রম করা (ôtikrôm kôrā)}]{\bn{অতিক্রম} \bn{করা} (ôtikrôm kôrā) — to surpass; to get past}

\label{lesson:BN-C377-otikrom-kora}



\begin{quote}
Before the new one: say the Bengali for progress, then the Bengali for speed.
\end{quote}

\subsection*{You'll want to know: \bn{অতিক্রম} \bn{করা}}
\textbf{\bn{অতিক্রম} \bn{করা}} — \emph{ôtikrôm kôrā} — \textquotedblleft{}to surpass; to get past\textquotedblright{}.

\begin{quote}
\bn{আমরা} \bn{যদি} \bn{এই} \bn{সমস্যা} \bn{অতিক্রম} \bn{করি}, \bn{তাহলে} \bn{কাজ} \bn{সহজ} \bn{হবে।}
\end{quote}

\begin{quote}
\emph{āmrā jôdi ei sômôsyā ôtikrôm kôri, tāhôle kāj sôhôj hôbe.}
\end{quote}

\begin{quote}
If we get past this problem, the work will be easier.
\end{quote}

Past (he/she …): \bn{তিনি} \bn{সব} \bn{সমস্যা} \bn{অতিক্রম} \bn{করলেন}, \emph{tini sôb sômôsyā ôtikrôm kôrlen}.

\begin{cousinweb}
\textbf{\bn{অতিক্রম} \bn{করা}} (\emph{ôtikrôm kôrā}) is Sanskrit \emph{atikrama}, passing over (Monier-Williams), with \bn{করা}: to go beyond an obstacle, a record or a limit.
\end{cousinweb}

\begin{grammarlens}[title={What you've built}]
One more word for saying what happens if.
\end{grammarlens}

\subsection*{Your turn}
\begin{itemize}
\raggedright
  \item \emph{Say it:} \emph{ôtikrôm kôrā}
  \item \emph{Say it:} \emph{ôtikrôm kôrā}, once more
  \item \emph{Say it:} \emph{āmrā jôdi ei sômôsyā ôtikrôm kôri, tāhôle kāj sôhôj hôbe.} — If we get past this problem, the work will be easier.
  \item \emph{Recall:} say the Bengali for progress, then the Bengali for speed, then say \emph{ôtikrôm kôrā} again
\end{itemize}

\begin{tcolorbox}[breakable,colback=teal!4,colframe=teal!35!black,title={Before you move on}]
What does \bn{অতিক্রম} \bn{করা} mean? (\textquotedblleft{}To surpass; to get past\textquotedblright{}.) Say it once more.
\end{tcolorbox}
\section[\texorpdfstring{\bn{সাধন} \bn{করা} (sādhôn kôrā)}{সাধন করা (sādhôn kôrā)}]{\bn{সাধন} \bn{করা} (sādhôn kôrā) — to accomplish; to bring about}

\label{lesson:BN-C377-sadhon-kora}



\begin{quote}
Before the new one: say the Bengali for speed, then the Bengali for to surpass.
\end{quote}

\subsection*{You'll want to know: \bn{সাধন} \bn{করা}}
\textbf{\bn{সাধন} \bn{করা}} — \emph{sādhôn kôrā} — \textquotedblleft{}to accomplish; to bring about\textquotedblright{}.

\begin{quote}
\bn{সবাই} \bn{একসঙ্গে} \bn{কাজ} \bn{করলে} \bn{আমরা} \bn{অনেক} \bn{কিছু} \bn{সাধন} \bn{করতে} \bn{পারব।}
\end{quote}

\begin{quote}
\emph{sôbāi eksôngge kāj kôrle āmrā ônek kichhu sādhôn kôrte pārbô.}
\end{quote}

\begin{quote}
If everyone works together, we can accomplish a lot.
\end{quote}

Past (he/she …): \bn{তিনি} \bn{তাঁর} \bn{লক্ষ্য} \bn{সাধন} \bn{করলেন}, \emph{tini tā̃r lôkkhyô sādhôn kôrlen}.

\begin{cousinweb}
\textbf{\bn{সাধন} \bn{করা}} (\emph{sādhôn kôrā}) is Sanskrit \emph{sādhana}, bringing about or accomplishing (Monier-Williams), with \bn{করা}.
\end{cousinweb}

\begin{grammarlens}[title={What you've built}]
One more word for saying what happens if.
\end{grammarlens}

\subsection*{Your turn}
\begin{itemize}
\raggedright
  \item \emph{Say it:} \emph{sādhôn kôrā}
  \item \emph{Say it:} \emph{sādhôn kôrā}, once more
  \item \emph{Say it:} \emph{sôbāi eksôngge kāj kôrle āmrā ônek kichhu sādhôn kôrte pārbô.} — If everyone works together, we can accomplish a lot.
  \item \emph{Recall:} say the Bengali for speed, then the Bengali for to surpass, then say \emph{sādhôn kôrā} again
\end{itemize}

\begin{tcolorbox}[breakable,colback=teal!4,colframe=teal!35!black,title={Before you move on}]
What does \bn{সাধন} \bn{করা} mean? (\textquotedblleft{}To accomplish; to bring about\textquotedblright{}.) Say it once more.
\end{tcolorbox}
